% ======================================================================
%  CAPÍTULO 7 — CONCLUSÕES
%
%  Roteiro: monografia/verificacao/anotacoes/cap07.md. Números das tabelas
%  de monografia/tabelas/ e da verificação (relatorio.md, seção 2e).
% ======================================================================
\chapter{Conclusões}
\label{cap:conclusoes}

Este trabalho projetou três controladores de atitude baseados em modelo ---
LQI, LQR com observador de estado estendido e LQG ---, embarcou-os no
firmware ArduPilot como \textit{backends} do \codigo{AC\_CustomControl} e os
comparou ao PID nativo, sintonizado pelo AutoTune, em SITL, sob as mesmas
referências e as mesmas perturbações reprodutíveis. Este capítulo responde à
pergunta de pesquisa, discute os limites que nenhuma lei de controle rompe e
as limitações do simulador, situa outras técnicas em relação a esses limites
e termina com uma recomendação para uma aeronave real e com sugestões de
trabalhos futuros.

% ----------------------------------------------------------------------
\section{Resposta à pergunta de pesquisa}
\label{sec:resposta}

A pergunta do Capítulo~\ref{cap:introducao} era se um controlador projetado
por métodos de controle moderno, a partir de um modelo da aeronave, é capaz
de igualar ou superar o PID nativo quando embarcado no firmware e submetido
às mesmas condições. A resposta depende do eixo e do controlador.

Em rolagem e arfagem, com a mesma frequência de cruzamento e a mesma margem
de fase do nativo, o LQI e o LQR com ESO \emph{não} o superam na rejeição de
torque: o pico do degrau curto de rolagem é de \ang{0.98} e \ang{0.99},
contra \ang{0.89} do nativo, e o do degrau sustentado, de \ang{0.65} e
\ang{0.66}, contra \ang{0.60}. Eles recuperam mais depressa ---
\SI{0.44}{\second} e \SI{0.42}{\second} no degrau sustentado, contra
\SI{2.20}{\second} ---, mas devolvem, quando um torque curto acaba, um
desvio de sinal oposto quase do tamanho do primeiro. E o LQI é, nesses
eixos, a própria cascata nativa com outros ganhos: o nativo voado com os
ganhos do LQI o reproduz (\ang{0.99} contra \ang{0.98} no degrau curto;
\ang{0.66} contra \ang{0.65} no degrau). O que o projeto muda é a sintonia, e
não a estrutura. No modelo, sob as mesmas restrições de ruído e de margem e
no espaço de ganhos varrido, nenhum PID em cascata supera o AutoTune no pico
do degrau curto de rolagem (Tabela~\ref{tab:projeto-portao}).

O LQG projetado para o ruído no esforço do nativo é o único que reduz o pico
em rolagem e arfagem: \ang{0.83} e \ang{0.82} no degrau curto (7\,\% e
11\,\% abaixo do nativo), \ang{0.55} no degrau sustentado (8\,\% e 11\,\%
abaixo), com recuperação de \SI{0.20}{\second}, e de 8\,\% a 9\,\% menos no
impulso de \SI{20}{\milli\second}. Com ruído de giroscópio e cinco voos por
par, a vantagem no impulso se confirma: $-9 \pm 2$\,\% em rolagem e
$-10 \pm 2$\,\% em arfagem. Ela se mantém com erro de modelo e com falha de
motor, cresce com o filtro do giroscópio de \SI{10}{\hertz} e o nativo
ressintonizado (\ang{1.72} contra \ang{2.13}, 19\,\% abaixo) e aparece também
sob torque estocástico, em que ele é o único abaixo do nativo
($-4 \pm 3$\,\% e $-6 \pm 2$\,\%), enquanto o LQI e o ESO ficam de 18\,\% a
23\,\% acima. O LQG com a regra de sintonia comum, de banda igual à do LQI,
não reduz o pico do degrau curto e do degrau, mas reduz o do impulso: a
vantagem vem da estrutura e do critério do portão juntos. A ablação mostra a
parte da estrutura: sem modelar o filtro do giroscópio ou o erro do ângulo
do AHRS, nenhum LQG cumpre o ruído no esforço do nativo, e o atraso de um
laço e o viés pesam pouco. A lei de realimentação é a mesma do LQR com ESO.

Em guinada, todos os propostos eliminam a recuperação lenta que a razão
$k_I = k_P/10$ do AutoTune impõe ao nativo: 0,75 a \SI{2.3}{\second}, contra
\SI{22}{\second}, no degrau sustentado. Mas o nativo com a guinada
ressintonizada também a elimina (\SI{1.8}{\second}), com o mesmo ruído no
esforço: a recuperação lenta é de sintonia. O LQI e o ESO também reduzem o
pico --- \ang{1.87} e \ang{1.78}, contra \ang{2.56} ---, mas com 2,4 e 2,0
vezes o ganho de ruído do nativo; com o ruído do nativo, o LQI piora para
\ang{6.16}. O LQG de mesmo ruído reduz o pico de 7\,\% a 8\,\% (\ang{2.35}
no degrau), com 1,13 vez o ruído no esforço do nativo, onde o melhor PID
reduz de 3\,\% a 4\,\%; em troca, recupera do degrau curto em
\SI{1.45}{\second}, contra \SI{0.35}{\second}, e fica 9\,\% acima sob torque
estocástico.

Por fim, a vantagem mais visível dos propostos, no rastreamento do degrau de
atitude --- sobressinal de 1\,\% a 6\,\%, contra 20\,\% a 27\,\%, em rolagem e
arfagem ---, vem da alimentação direta pela inversão do modelo, e não da
realimentação: o PID com a mesma alimentação direta chega a 4,3\,\% de
sobressinal em rolagem, contra 5,8\,\% do LQI.

Em resumo, um controlador baseado em modelo, embarcado no firmware, iguala o
PID nativo na rejeição de perturbações em rolagem e arfagem e só o supera
quando modela o caminho de medida e é projetado sob a mesma restrição de
ruído (o LQG, de 7\,\% a 11\,\% no pico). Em guinada, a recuperação rápida
vem da sintonia, e a redução do pico vem de banda e de ruído ou, no LQG, é
de alguns por cento. E ele supera o nativo no rastreamento por uma
alimentação direta que o próprio PID poderia usar.

% ----------------------------------------------------------------------
\section{Limites que nenhuma lei de controle rompe}
\label{sec:limites}

Os resultados mostram limites que vêm da informação disponível ao controlador
e da planta, e não da lei de controle.

\subsection{A integral de Bode}

Para qualquer controlador linear e invariante
no tempo, atenuar a perturbação numa faixa de frequências aumenta a
sensibilidade em outra (Seção~\ref{sec:bode}). O que as varreduras medem é
$P\,S$, e não $S$; como a planta $P$ é a mesma para todos os controladores, a
diferença entre os $\ln|P\,S|$ de dois deles é a diferença entre os seus
$\ln|S|$, e o argumento vale para a comparação. No modelo, a integral
discreta de $\ln|S|$ é nula em todos os controladores, com as áreas negativa
e positiva iguais~\cite{sung1988,mohtadi1990}; o LQG de mesmo ruído tem a
maior área em rolagem, 0,051 contra 0,049 do nativo: atenua mais em baixa
frequência, o que dá a recuperação rápida, e paga com um máximo de $|P\,S|$
maior perto do cruzamento. No tempo, a mesma troca aparece como o desvio de
sinal oposto que os propostos devolvem quando um torque curto acaba.

\subsection{O ruído do giroscópio contra a autoridade em guinada}

A
eficácia de controle em guinada é de \SI{14.37}{\radian\per\second\squared},
30 vezes menor que em rolagem (Tabela~\ref{tab:modelo-linear}). Mais banda
em guinada custa ruído no esforço na mesma proporção, como mostram as
fronteiras de projeto e medida (Figuras~\ref{fig:projeto-fronteira}
e~\ref{fig:comparacao-fronteira}).

\subsection{O filtro do giroscópio e o atraso do laço}

O filtro de
\SI{20}{\hertz} custa cerca de \ang{33} de fase no cruzamento da rolagem, em
\SI{48.5}{\radian\per\second}, e o ângulo do AHRS chega com um laço de atraso,
cerca de \ang{7} na mesma frequência. A margem de fase da malha ideal do
LQR, de 65 a \ang{73}, cai por isso para 39 a \ang{42} medidos no SITL;
elas não sobrevivem ao sensor (Seção~\ref{sec:projeto-margens}).

O LQG é ótimo para o critério quadrático e para as estatísticas de ruído e
de perturbação assumidas no projeto --- não para as métricas dos ensaios ---,
e usa só o ângulo do AHRS e o giroscópio filtrado. Ganhos maiores viriam de
mais informação e da cadeia de sinal e de atuação, e beneficiariam também o
PID: um
filtro \textit{notch} dinâmico no lugar do passa-baixas; a telemetria de rotação dos controladores dos motores; ou um
giroscópio menos ruidoso ou mais bem isolado da vibração.

% ----------------------------------------------------------------------
\section{Limitações do simulador e do estudo}
\label{sec:limitacoes}

O SITL é um cenário otimista para o controle baseado em modelo:
\begin{alineas}
    \item não tem dinâmica de motor, só um limitador de taxa do comando
          (\codigo{SIM\_FRM\_SLEW\_MAX}), que não atua na regulação
          (Seção~\ref{sec:motor-sitl}); numa aeronave real, cada motor
          acrescentaria um polo entre o comando e o empuxo;
    \item não modela o termo giroscópico
          $\vetor{\omega}\times\matriz{J}\,\vetor{\omega}$ nem o efeito
          giroscópico dos rotores (Seção~\ref{sec:sitl}), e por isso o
          acoplamento entre os eixos foi desprezível;
    \item aplica a carga suspensa e o cabo (\codigo{SIM\_SLUP\_*} e
          \codigo{SIM\_TETH\_*}) só como força sobre o centro de massa, sem
          momento;
    \item trata o vento quase só como perturbação de posição: o arrasto de
          momento dos rotores só produz momento quando os empuxos diferem, e
          uma rajada de \SI{10}{\metre\per\second} desviou a rolagem em
          \ang{0.015} a \ang{0.018} (Seção~\ref{sec:res-vento});
    \item nos ensaios-base, entrega um giroscópio praticamente sem ruído,
          só com o ruído mínimo do simulador (Seção~\ref{sec:sensores-sitl});
    \item tem a dinâmica angular quase independente da massa, porque a
          propulsão é refeita a partir dela e a inércia escala junto; o erro
          de modelo dos ensaios foi criado pelo empuxo de pairado e pelo
          tamanho do quadro.
\end{alineas}

O estudo tem também limitações próprias. Com ruído de giroscópio, voos
idênticos variam até 8\,\%, e só o impulso com
\codigo{SIM\_GYR\_RND} = 20 voa cinco vezes por par; nas tabelas de níveis
de ruído, com um voo por par, diferenças de poucos por cento entre
controladores não são conclusivas (Seção~\ref{sec:res-ruido}). As métricas
estão na atitude do AHRS, que tem um erro lento de rumo de até cerca de
\ang{1}, igual com qualquer controlador e invisível a todos; a atitude
verdadeira está no Apêndice~\ref{ap:bases}. No LQG de mesmo ruído, o ruído
no esforço medido em guinada ficou 13\,\% acima do previsto, e, depois do
degrau sustentado de guinada, resta um erro de cerca de \ang{0.03} que decai
em dezenas de segundos; as duas coisas pedem investigação antes de um voo
real.

% ----------------------------------------------------------------------
\section{Outras técnicas de controle}
\label{sec:outras-tecnicas}

Os limites da Seção~\ref{sec:limites} valem para toda lei linear e invariante
no tempo com a mesma informação. Outras técnicas atacariam pontos diferentes,
e cada uma teria ganho num grupo específico dos ensaios deste trabalho.

\subsection{Inversão dinâmica não linear incremental (INDI)}

A INDI controla
o incremento de esforço a partir da aceleração angular medida, e com isso
depende pouco do modelo da aeronave e absorve momentos não
modelados~\cite{sieberling2010}; foi aplicada à atitude de multirrotores com
estimação adaptativa da eficácia de controle~\cite{smeur2016}. Ganharia nos
ensaios de erro de modelo e de falha de motor, sobretudo com a realimentação
da rotação dos motores. Em rejeição pura no SITL, ficaria próxima do ESO e do
LQG, que também estimam a perturbação, e, em guinada, limitada pelo ruído do
mesmo modo, porque depende da derivada do giroscópio. É a alternativa mais
relevante para uma aeronave real.

\subsection{\texorpdfstring{Controle $H_\infty$ e \textit{loop shaping}}{Controle H-infinito e loop shaping}}

O projeto
$H_\infty$ trata o pior caso da sensibilidade~\cite{zames1981}, e o
\textit{loop shaping} $H_\infty$ dá margens de robustez por
construção~\cite{mcfarlane1992}; foi aplicado à atitude de um quadricóptero
real sobre um modelo identificado com incerteza
multiplicativa~\cite{noormohammadiasl2020}. Teria ganho onde o critério
quadrático do LQG não olha: poderia limitar por construção o máximo de
$|P\,S|$ do LQG perto do cruzamento e garantiria margens, que o LQG não
tem~\cite{doyle1978}, provavelmente com algum custo no pico do degrau
curto.

\subsection{Controle adaptativo}

O controle adaptativo por modelo de
referência~\cite{astrom1995} e o controle adaptativo
$\mathcal{L}_1$~\cite{hovakimyan2010} ajustam o controlador em voo; o
primeiro foi avaliado em voo num quadricóptero com perda de
empuxo~\cite{dydek2013}, e o segundo, como complemento de um controle
geométrico~\cite{wu2022}. Ganhariam quando a eficácia de controle varia ---
com a bateria, a carga ou o desgaste ---, o que, no SITL, aparece só nos
ensaios de erro de modelo e de falha de motor.

\subsection{Controle preditivo (MPC)}

O MPC trata restrições
explicitamente~\cite{mayne2000} e foi aplicado ao rastreamento de trajetória
de multirrotores, com vento~\cite{kamel2017}. Sem restrições ativas, a lei
coincide com a do LQR, e o ganho ficaria nos ensaios de saturação e de
ângulos grandes, em que a folga de guinada da alocação e a saturação dos
motores pesam, e na antecipação do alvo. A \SI{400}{\hertz}, porém, o custo
de calcular a lei a cada laço é alto para a controladora de voo.

\subsection{Modos deslizantes}

O controle por modos
deslizantes~\cite{utkin1977} e o algoritmo \textit{super-twisting}, de
segunda ordem~\cite{levant1993}, rejeitam em tempo finito as perturbações
que entram no canal do controle e foram aplicados à atitude de um
quadricóptero real~\cite{derafa2012}. São, porém, sensíveis ao ruído de
medida: seriam piores que os propostos em guinada, onde o ruído já limita a
banda, e produzem oscilação de alta frequência no esforço
(\textit{chattering}).

\subsection{\texorpdfstring{Controle geométrico em $SO(3)$}{Controle geométrico em SO(3)}}

O controle formulado no grupo das
rotações~\cite{bullo1999}, aplicado ao quadricóptero em
$SE(3)$~\cite{lee2010}, evita as singularidades e as aproximações de ângulos
pequenos. Perto do pairado, equivale a um controle proporcional-derivativo, e
só teria ganho em manobras grandes, como as dos ensaios de \ang{30} e
\ang{45} (Tabela~\ref{tab:res-grande}).

% ----------------------------------------------------------------------
\section{Recomendação para uma aeronave real}
\label{sec:recomendacao}

Os resultados permitem uma recomendação de engenharia, que é uma opinião
baseada nos ensaios em SITL, e não um resultado deles. Para melhorar o
controle de atitude de uma aeronave real com o ArduPilot, nesta ordem:
\begin{alineas}
    \item acrescentar ao PID nativo a alimentação direta pela inversão do
          modelo, aplicada ao alvo modelado e somada ao esforço, como no
          \textit{backend} de exemplo (\codigo{cc\_pid\_ff}), que dá quase
          toda a vantagem dos propostos no rastreamento
          (Tabela~\ref{tab:res-ff}). Os termos de alimentação direta do
          próprio PID de taxa do ArduPilot (\codigo{FF} e \codigo{D\_FF})
          não fazem o mesmo: agem sobre o alvo de velocidade, que inclui a
          realimentação do ângulo, e mudam a malha. No modelo, eles subiriam
          o cruzamento de rolagem de 46,6 para \SI{68}{\radian\per\second},
          com a margem de fase de \ang{38.8} para \ang{35.4} e 1,6 vez o
          ruído no esforço; em guinada, o \codigo{D\_FF} necessário fica fora
          da faixa documentada. Por isso, não foram voados;
    \item reforçar a ação integral em guinada: o nativo com a guinada
          ressintonizada recupera do degrau de torque em \SI{1.8}{\second},
          contra \SI{22}{\second}, sem mais ruído no esforço, mas quase não
          reduz o pico (Tabela~\ref{tab:res-nativo-guinada}); um observador
          de estado estendido reduz o pico com mais banda e mais ruído;
    \item considerar o LQG só depois de identificar a dinâmica dos motores e
          o ruído real de vibração na aeronave, porque o SITL é um cenário
          otimista para o controle baseado em modelo e a vantagem dele, de
          7\,\% a 11\,\% no pico, é pequena diante dessas incertezas; e
          investigar antes o ruído no esforço em guinada, 13\,\% acima do
          previsto.
\end{alineas}

% ----------------------------------------------------------------------
\section{Trabalhos futuros}
\label{sec:trabalhos-futuros}

Ficam como continuidade natural deste trabalho:
\begin{alineas}
    \item voar os controladores numa aeronave real, começando pela
          alimentação direta no PID nativo;
    \item investigar o ruído no esforço do LQG em guinada e o erro residual
          que ele deixa depois do degrau sustentado;
    \item repetir os ensaios de níveis de ruído de giroscópio com vários voos
          por par, como já se fez no impulso com
          \codigo{SIM\_GYR\_RND} = 20;
    \item modelar o erro lento do AHRS e descontá-lo da atitude verdadeira,
          para comparar os controladores na grandeza física sem esse erro;
    \item avaliar o filtro \textit{notch} dinâmico e a INDI no mesmo
          procedimento de ensaio;
    \item estender o procedimento às malhas externas e aos ensaios de vento
          com momento, num simulador que modele a aerodinâmica dos rotores.
\end{alineas}
